\documentclass[11pt]{article}
\usepackage{preamble}

\title{Matsumoto's theorem with coefficients}
\author{Working note for \texttt{mainv2}}
\date{\today}

\newcommand{\EN}{\mathcal{E}_N}
\newcommand{\Coc}{\mathrm{Coc}}
\newcommand{\Fam}{\mathrm{Fam}}
\newcommand{\altw}[3]{\Pi(#1,#2:#3)}

\begin{document}
\maketitle

\noindent
This note proves Theorem~1.1 of \texttt{mainv2} by a direct route. The Coxeter input is Tits's
word property and two facts about minimal coset representatives, all cited. The proof is stated for
an arbitrary presheaf on the category $\EN$ of Paliga and Ziemia\'nski; precubical sets enter only
in \Cref{sec:cubes}. Locators for \cite{Paris} and \cite{PZ} refer to their arXiv v1.

\section{Coxeter background}

Let $S_N$ be the symmetric group on $[N]$, $s_t := (t,\,t+1)$ for $t \in [N-1]$, and $S :=
\{s_1,\dots,s_{N-1}\}$, ordered by index. Then $(S_N,S)$ is a Coxeter system \cite[Prop.~1.5.4]{BB},
with $m_{s_i,s_k} = \cox(i,k)$, which is $3$ if $|i-k| = 1$ and $2$ otherwise \cite[(1.29)]{BB}. Write
$\ell$ for its length. For $X \subseteq S$, $W_X$ is the subgroup generated by $X$; an element $w$ is
\textbf{$(\emptyset,X)$-minimal} if it is of minimal length in $wW_X$ \cite[\S2.2]{Paris}, and $W^X$
is the set of these. For $s,t \in S$ and $k \in \N$, $\altw stk$ is the alternating word $sts\cdots$
of length $k$, extending the notation of \cite[\S1]{Paris} (defined there for $k \ge 2$) to $k = 0,1$,
and also the element of $S_N$ it represents. The next two results are stated in \cite{Paris} for an
arbitrary Coxeter system $(W,S)$; we use them for $(S_N,S)$.

\begin{theorem}[{\cite[Thm.~2.7]{Paris}}]\label{thm:tits}
  Let $w \in W$, and let $\mu, \mu'$ be two reduced expressions of $w$. Then there is a finite
  sequence of elementary $M$-transformations joining $\mu$ to $\mu'$.
\end{theorem}

Here an \textbf{elementary $M$-transformation} joins $\mu$ to $\mu'$ when there are words
$\nu_1,\nu_2$ and $s,t \in S$ with $m_{s,t} \neq \infty$, $\mu = \nu_1\,\altw st{m_{s,t}}\,\nu_2$ and
$\mu' = \nu_1\,\altw ts{m_{s,t}}\,\nu_2$ \cite[\S2.2]{Paris}.

\begin{proposition}[{\cite[Prop.~2.9]{Paris}}]\label{prop:min}
  Let $(W,S)$ be a Coxeter system.
  \begin{enumerate}
    \item[(2)] Let $X \subset S$, and let $w \in W$. Then $w$ is $(\emptyset,X)$-minimal if and
      only if $\ell(ws) > \ell(w)$ for all $s \in X$, and $\ell(ws) > \ell(w)$ for all $s \in X$ if
      and only if $\ell(wu) = \ell(w) + \ell(u)$ for all $u \in W_X$.
    \item[(4)] Let $X \subset S$, and let $w \in W_X$. If $\mu = s_1 \cdots s_\ell$ is a reduced
      expression of $w$, then $s_1, \dots, s_\ell \in X$.
  \end{enumerate}
\end{proposition}

\begin{proposition}\label{prop:typeA}
  \begin{enumerate}
    \item \cite[Prop.~1.4.2(ii)]{BB} $\ell(uw) \equiv \ell(u) + \ell(w) \pmod 2$ for all $u, w$.
    \item \cite[Prop.~1.5.2]{BB} $\ell(x) = \mathrm{inv}(x)$, where $\mathrm{inv}(x) := \#\{(i,j) :
      i < j,\ x(i) > x(j)\}$ \cite[(1.25)]{BB}.
    \item \cite[Prop.~1.5.3]{BB} $D_R(x) = \{s_i \in S : x(i) > x(i+1)\}$, where $D_R(x) := \{s \in
      S : \ell(xs) < \ell(x)\}$ \cite[(1.19), (1.20)]{BB}.
  \end{enumerate}
\end{proposition}

\begin{lemma}\label{lem:ascent}
  For $x \in S_N$ and $t \in [N-1]$: $\ell(xs_t) > \ell(x)$ if and only if $x(t) < x(t+1)$.
\end{lemma}
\begin{proof}
  $\ell(xs_t) \neq \ell(x)$ by \Cref{prop:typeA}(1), so $\ell(xs_t) > \ell(x)$ if and only if $s_t
  \notin D_R(x)$, which is $x(t) < x(t+1)$ by \Cref{prop:typeA}(3).
\end{proof}

\begin{lemma}\label{lem:prefix}
  If $x_1 \cdots x_k$ is a reduced expression, then so is $x_1 \cdots x_i$ for every $i \le k$.
\end{lemma}
\begin{proof}
  If $x_1\cdots x_i$ had an expression of length $< i$, appending $x_{i+1}\cdots x_k$ would give an
  expression of $x_1\cdots x_k$ of length $< k$.
\end{proof}

\begin{lemma}\label{lem:dihedral}
  Let $s \neq t$ in $S$ and $m := m_{s,t}$. Then $\altw stk$ is a reduced expression for every $k \le
  m$. Consequently, if $x$ is the $(j+1)$-st letter of $\altw stm$, then $\altw stj \in W^{\{x\}}$ for
  every $j < m$.
\end{lemma}
\begin{proof}
  By \Cref{lem:prefix} it suffices to take $k = m$, and by \Cref{prop:typeA}(2) to show
  $\mathrm{inv}(\altw stm) = m$. If $m = 2$, $\altw st2$ is a product of two transpositions of disjoint
  pairs of adjacent positions, so it has $2$ inversions. If $m = 3$, $\{s,t\} = \{s_i, s_{i+1}\}$ and
  $\altw st3$ is the transposition $(i,\,i+2)$, which reverses the three positions $i, i+1, i+2$ and
  fixes the rest, so it has $3$ inversions. For the last sentence, $\altw stj\,x = \altw st{j+1}$, so
  $\ell(\altw stj\,x) = j+1 > \ell(\altw stj)$, and \Cref{prop:min}(2) applies with $X = \{x\}$.
\end{proof}

\section{The category \texorpdfstring{$\EN$}{E\_N} and its presheaves}

\begin{definition}[{\cite[Def.~6.11]{PZ}}]\label{def:EN}
  Let $\EN$ be the category whose objects are subsets of $\{1,\dots,N-1\}$. For $B = \{b_1 < \cdots <
  b_{r-1}\} \supseteq B' = \{b'_1 < \cdots < b'_{r'-1}\}$, morphisms from $B$ to $B'$ are permutations
  $\phi$ of $\{1,\dots,N\}$ satisfying
  \begin{enumerate}
    \item[(x)] $\phi(\{b'_{p-1}+1,\dots,b'_p\}) = \{b'_{p-1}+1,\dots,b'_p\}$ for $p \in \{1,\dots,r'\}$
      (with $b'_0 = 0$, $b'_{r'} = N$);
    \item[(y)] $\phi(b_{p-1}+1) < \phi(b_{p-1}+2) < \cdots < \phi(b_p)$ for $p \in \{1,\dots,r\}$ (with
      $b_0 = 0$, $b_r = N$).
  \end{enumerate}
  For $B \not\supseteq B'$, $\EN(B,B') = \emptyset$. Composition is composition of permutations.
\end{definition}

\cite[Def.~6.11]{PZ} also notes that (x) is equivalent to
\begin{enumerate}
  \item[(x$'$)] for all $i,j$: if $i \le b' < j$ for some $b' \in B'$, then $\phi(i) < \phi(j)$.
\end{enumerate}
For $J \subseteq S$ put $B(J) := \{t \in [N-1] : s_t \notin J\}$. Then $t$ and $t+1$ lie in one block
$\{b_{p-1}+1,\dots,b_p\}$ of $B(J)$ exactly when $s_t \in J$.

\begin{lemma}\label{lem:EN}
  Let $J' \subseteq J \subseteq S$. Then $\EN(B(J'),B(J)) = W_J \cap W^{J'}$.
\end{lemma}
\begin{proof}
  \emph{(y) for $B(J')$ is $\phi \in W^{J'}$.} (y) says $\phi(t) < \phi(t+1)$ whenever $t, t+1$ lie in one
  block of $B(J')$, that is whenever $s_t \in J'$. By \Cref{lem:ascent} this is $\ell(\phi s) >
  \ell(\phi)$ for all $s \in J'$, which is $\phi \in W^{J'}$ by \Cref{prop:min}(2).

  \emph{(x) for $B(J)$ is $\phi \in W_J$.} Each $s_t \in J$ exchanges two positions of one block of
  $B(J)$, so it satisfies (x), and so does every element of $W_J$. Conversely let $\phi$ satisfy (x);
  we show $\phi \in W_J$ by induction on $\ell(\phi)$. If $\phi \neq 1$ then $\phi(t) > \phi(t+1)$ for
  some $t$, so $s_t \in D_R(\phi)$ (\Cref{prop:typeA}(3)). By (x$'$), $t \notin B(J)$, that is $s_t \in
  J$. Then $\phi s_t$ satisfies (x) and $\ell(\phi s_t) < \ell(\phi)$, so $\phi s_t \in W_J$ by
  induction, and $\phi \in W_J$.
\end{proof}

From now on we write $J$ for the object $B(J)$ of $\EN$, so $\EN(J',J) = W_J \cap W^{J'}$ for $J'
\subseteq J$ and $\emptyset$ otherwise, composed by the product of $S_N$. In particular $\EN(\emptyset,J)
= W_J$, since (y) is empty for $B(\emptyset) = [N-1]$.

Fix a presheaf $P$ on $\EN$. For $v \in \EN(J',J)$ and $d \in P(J)$ write $d \cdot v \in P(J')$, so
$(d\cdot v)\cdot u = d \cdot (vu)$ and $d \cdot 1 = d$. The elements of $P(\emptyset)$ are the
\textbf{runs}. Every $v \in \EN(J',J)$ also lies in $\EN(\emptyset,J) = W_J$, and $(d\cdot v)\cdot 1
= d\cdot v$ as runs; we write $d\cdot v$ in both cases.

The category of elements $\int P$ has objects the $d \in P(J)$, $J \subseteq S$, and a morphism $c \to
d$, for $c \in P(J')$ and $d \in P(J)$, is a $v \in \EN(J',J)$ with $d\cdot v = c$; composition is
the product. Let $W_P$ be the class of morphisms $v = 1$ of $\int P$, and
\[
  \Br(P) \;:=\; (\textstyle\int P)\op\bigl[(W_P\op)\inv\bigr],
\]
with localization functor $Q$, contravariant on $\int P$ as in \texttt{mainv2}. For $d \in P(J)$
and $w \in W_J$, write $\rho^d_w : d\cdot w \to d$ for the morphism $w$ of $\int P$. Then
$\rho^d_1 \in W_P$.

\begin{definition}\label{def:art}
  The \textbf{Artin polygraph} $\Art(P)$ has the runs as $0$-generators. A $1$-generator from $X$ to
  $Y$ is an $\alpha \in P(\{s\})$, $s \in S$, with $\alpha\cdot1 = X$ and $\alpha\cdot s = Y$. A
  $2$-generator is an $\varepsilon \in P(\{s,t\})$ with $s < t$, and, for $m := m_{s,t}$,
  \[
    \src(\varepsilon) = \bigl(\varepsilon\cdot\altw st0,\ \dots,\ \varepsilon\cdot\altw st{m-1}\bigr),
    \qquad
    \tgt(\varepsilon) = \bigl(\varepsilon\cdot\altw ts0,\ \dots,\ \varepsilon\cdot\altw ts{m-1}\bigr).
  \]
\end{definition}

By \Cref{lem:dihedral}, $\altw stj \in \EN(\{x\},\{s,t\})$ for the $(j+1)$-st letter $x$ of $\altw
stm$, so $\varepsilon\cdot\altw stj$ is a $1$-generator from $\varepsilon\cdot\altw stj$ to
$\varepsilon\cdot\altw st{j+1}$, and $\src(\varepsilon)$ is a path from $\varepsilon\cdot 1$ to
$\varepsilon\cdot\altw stm$. Likewise for $\tgt(\varepsilon)$, which ends at the same run since $\altw
stm = \altw tsm$ in $S_N$ \cite[(1.29)]{BB}.

\section{Cocycles and Artin families}

Let $D$ be a category.

\begin{definition}\label{def:cocycle}
  A \textbf{$P$-cocycle} in $D$ is an object $x_X$ of $D$ for each run $X$ and an arrow
  $\sigma_d(w) : x_{d\cdot1} \to x_{d\cdot w}$ for each $J \subseteq S$, $d \in P(J)$ and $w \in
  W_J$, such that
  \begin{enumerate}[label=(C\arabic*)]
    \item $\sigma_d(1) = \id$;
    \item $\sigma_d(vu) = \sigma_{d\cdot v}(u) \circ \sigma_d(v)$ for all $J' \subseteq J$,
      $v \in \EN(J',J)$ and $u \in W_{J'}$.
  \end{enumerate}
\end{definition}

\begin{definition}\label{def:family}
  A \textbf{$P$-Artin family} in $D$ is an object $x_X$ of $D$ for each run $X$ and an arrow $a_\alpha :
  x_{\alpha\cdot1} \to x_{\alpha\cdot s}$ for each $\alpha \in P(\{s\})$, $s \in S$, such that for
  every $2$-generator $\varepsilon \in P(\{s,t\})$ of $\Art(P)$, with $m = m_{s,t}$,
  \[
    a_{\varepsilon\cdot\altw st{m-1}} \circ \cdots \circ a_{\varepsilon\cdot\altw st0}
    \;=\;
    a_{\varepsilon\cdot\altw ts{m-1}} \circ \cdots \circ a_{\varepsilon\cdot\altw ts0} .
  \]
\end{definition}

Write $\Coc_P(D)$ and $\Fam_P(D)$ for the sets of these. A functor $F : D \to D'$ acts on both by
applying $F$ to every object and arrow. Write $q : \Art(P)_1^* \to \pres{\Art(P)}$ for the quotient of
the free category on the $1$-generators.

\begin{proposition}\label{prop:pres}
  Sending a functor $F : \pres{\Art(P)} \to D$ to $(F(X), F(q\alpha))$ is a bijection onto
  $\Fam_P(D)$.
\end{proposition}
\begin{proof}
  The arrows $F(q\alpha)$ satisfy the family relations because $q(\src\varepsilon) =
  q(\tgt\varepsilon)$. Conversely a family extends uniquely to the free category
  \cite[Lem.~2.1.3]{Polygraphs}, and then uniquely to $\pres{\Art(P)}$ \cite[Lem.~2.3.4]{Polygraphs},
  since its relations say exactly that the source and target of each $2$-generator have equal image.
\end{proof}

\begin{theorem}[Matsumoto with coefficients]\label{thm:matsumoto}
  Sending $(x,\sigma)$ to $(x,a)$ with $a_\alpha := \sigma_\alpha(s)$, for $\alpha \in P(\{s\})$, is a
  bijection $\Coc_P(D) \to \Fam_P(D)$, compatible with the action of functors $D \to D'$.
\end{theorem}
\begin{proof}
  \emph{Well defined.} $\sigma_\alpha(s) : x_{\alpha\cdot1} \to x_{\alpha\cdot s}$ has the right type.
  Let $\varepsilon \in P(\{s,t\})$ be a $2$-generator and $x_{j+1}$ the $(j+1)$-st letter of $\altw
  stm$. By \Cref{lem:dihedral}, $\altw stj \in \EN(\{x_{j+1}\},\{s,t\})$, so (C2) with $v = \altw
  stj$ and $u = x_{j+1}$ gives
  \[
    \sigma_\varepsilon\bigl(\altw st{j+1}\bigr)
      = a_{\varepsilon\cdot\altw stj} \circ \sigma_\varepsilon\bigl(\altw stj\bigr) .
  \]
  By induction and (C1), the left side of the family relation is $\sigma_\varepsilon(\altw stm)$.
  Likewise the right side is $\sigma_\varepsilon(\altw tsm)$, and $\altw stm = \altw tsm$.

  \emph{The inverse.} Let $(x,a) \in \Fam_P(D)$, $d \in P(J)$ and $w \in W_J$. For a reduced
  expression $\mu = x_1\cdots x_k$ of $w$, put $w_i := x_1\cdots x_i$ and
  \[
    \sigma_d(\mu) := a_{d\cdot w_{k-1}} \circ \cdots \circ a_{d\cdot w_0} .
  \]
  This is well typed: $x_i \in J$ by \Cref{prop:min}(4); $w_{i-1}x_i = w_i$ is reduced by
  \Cref{lem:prefix}, so $\ell(w_{i-1}x_i) > \ell(w_{i-1})$ and $w_{i-1} \in W^{\{x_i\}}$ by
  \Cref{prop:min}(2); so $w_{i-1} \in \EN(\{x_i\},J)$, and $d\cdot w_{i-1}$ is a $1$-generator from
  $d\cdot w_{i-1}$ to $d\cdot w_i$.

  It does not depend on $\mu$. By \Cref{thm:tits} it suffices to compare $\mu = \nu_1\altw
  stm\nu_2$ with $\mu' = \nu_1\altw tsm\nu_2$, both reduced, and we may take $s < t$. Let $v$ be the
  element represented by $\nu_1$. Then $v \in W_J$, as the letters of $\mu$ lie in $J$
  (\Cref{prop:min}(4)), and $vs$, $vt$ are represented by prefixes of $\mu$, $\mu'$, so $\ell(vs) =
  \ell(vt) = \ell(v)+1$ by \Cref{lem:prefix}; hence $v \in W^{\{s,t\}}$ by \Cref{prop:min}(2), and
  $\varepsilon := d\cdot v \in P(\{s,t\})$ is a $2$-generator. The two composites have the same
  factors along $\nu_1$, and along $\nu_2$ since $\nu_1\altw stm$ and $\nu_1\altw tsm$ represent the
  same element. Along the middle, the factors are $a_{d\cdot v\altw stj} = a_{\varepsilon\cdot\altw
  stj}$ and $a_{\varepsilon\cdot\altw tsj}$, whose composites agree by the family relation of
  $\varepsilon$. Write $\sigma_d(w)$ for the common value.

  (C1) holds as the empty product. For (C2), let $v \in \EN(J',J)$ and $u \in W_{J'}$. By
  \Cref{prop:min}(2), $\ell(vu) = \ell(v) + \ell(u)$, so reduced expressions of $v$ and $u$
  concatenate to a reduced expression of $vu$. Reading $\sigma_d(vu)$ along it, the factors past $v$
  are $a_{d\cdot vu_i} = a_{(d\cdot v)\cdot u_i}$, which compose to $\sigma_{d\cdot v}(u)$.

  \emph{Mutually inverse.} Starting from a family, $\sigma_\alpha(s)$ is read along the expression
  $s$, so it is $a_{\alpha\cdot1} = a_\alpha$. Starting from a cocycle, (C2) with $v = w_{i-1}$ and
  $u = x_i$ gives $\sigma_d(w_i) = \sigma_{d\cdot w_{i-1}}(x_i) \circ \sigma_d(w_{i-1})$, so by
  induction $\sigma_d(w)$ is the composite of the $\sigma_{d\cdot w_{i-1}}(x_i) = a_{d\cdot
  w_{i-1}}$.

  Compatibility with functors is clear from the formulas.
\end{proof}

\section{The presentation}

\begin{lemma}\label{lem:s}
  Putting $x_X := Q(X)$ and $s_d(w) := Q(\rho^d_w) \circ Q(\rho^d_1)\inv$ defines a $P$-cocycle $s$ in
  $\Br(P)$.
\end{lemma}
\begin{proof}
  (C1) is clear. For (C2), let $v \in \EN(J',J)$, $u \in W_{J'}$ and $c := d\cdot v$, and let $g :
  c \to d$ be the morphism $v$. In $\int P$, $\rho^d_{vu} = g \circ \rho^c_u$ and $\rho^d_v = g
  \circ \rho^c_1$, as both sides are the same element with the same target. Since $Q$ reverses
  composition,
  \[
    s_d(vu) = Q(\rho^c_u)\,Q(g)\,Q(\rho^d_1)\inv
      = \bigl[Q(\rho^c_u)\,Q(\rho^c_1)\inv\bigr]\bigl[Q(\rho^c_1)\,Q(g)\,Q(\rho^d_1)\inv\bigr]
      = s_c(u) \circ s_d(v) . \qedhere
  \]
\end{proof}

\begin{theorem}\label{thm:A}
  Let $E : \pres{\Art(P)} \to \Br(P)$ be the functor that \Cref{prop:pres} assigns to the
  family of the cocycle $s$ (\Cref{thm:matsumoto}); so $E(X) = Q(X)$ and $E(q\alpha) =
  Q(\rho^\alpha_s)\,Q(\rho^\alpha_1)\inv$. Then $E$ is an equivalence of categories. For a map $\varphi
  : P \to P'$ of presheaves, $E_{P'} \circ \pres{\Art(\varphi)} = \Br(\varphi) \circ E_P$.
\end{theorem}
\begin{proof}
  Let $\Gamma$ be the cocycle in $\pres{\Art(P)}$ whose family is $(X, q\alpha)$
  (\Cref{thm:matsumoto}). Define $\Theta : (\int P)\op \to \pres{\Art(P)}$ by $\Theta(c) := c\cdot1$
  and $\Theta(g) := \Gamma_d(v)$ for $g = v : c \to d$; the types match because $c\cdot 1 = d\cdot
  v$. It is a functor by (C2), since composition in $\int P$ is the product, and it sends $W_P$ to
  identities by (C1). By \cite[\S I.1.1, Lem.~I.1.2]{GZ} there is a unique $\Phi : \Br(P) \to
  \pres{\Art(P)}$ with $\Phi Q = \Theta$.

  \emph{$\Phi E = \id$.} $\Phi E(X) = X\cdot1 = X$, and $\Phi E(q\alpha) = \Gamma_\alpha(s) \circ
  \Gamma_\alpha(1)\inv = q\alpha$. So $\Phi E$ and $\id$ agree on generators, and they are equal by
  the uniqueness in \cite[Lem.~2.1.3, Lem.~2.3.4]{Polygraphs}.

  \emph{$E\Phi \cong \id$.} $E\circ\Gamma$ is a cocycle in $\Br(P)$ with the same family as $s$, so
  $E\circ\Gamma = s$ by \Cref{thm:matsumoto}. Hence for $g = v : c \to d$, using $\rho^d_v = g
  \circ \rho^c_1$,
  \[
    E\Theta(g) = s_d(v) = Q(\rho^c_1)\,Q(g)\,Q(\rho^d_1)\inv .
  \]
  So the arrows $\eta_c := Q(\rho^c_1) : Q(c) \to E\Theta(c)$ form a natural isomorphism $Q \cong
  E\Theta = E\Phi Q$. Since precomposition with $Q$ is an isomorphism onto a full subcategory
  \cite[Lem.~I.1.2]{GZ}, $\eta = \eta' Q$ for a unique $\eta' : \id \Rightarrow E\Phi$, invertible
  because $\eta\inv$ descends too.

  \emph{Naturality.} $\int\varphi$ carries $W_P$ into $W_{P'}$ and $\rho^d_w$ to
  $\rho^{\varphi d}_w$, and $\Art(\varphi)$, defined by $\varphi$ on generators, is a morphism of
  polygraphs because $\varphi$ commutes with restriction. Both composites send $q\alpha$ to
  $Q(\rho^{\varphi\alpha}_s)\,Q(\rho^{\varphi\alpha}_1)\inv$, so they agree on generators.
\end{proof}

\begin{corollary}\label{cor:artin}
  At the terminal presheaf $P = 1$, $\Art(1)$ is the component at $N$ of Artin's braid
  polygraph
  (\texttt{def:artin-poly}).
\end{corollary}
\begin{proof}
  Each $P(J)$ is a point. So there is one $0$-generator, one $1$-generator $s_t$ for each $t$, and
  one $2$-generator for each $s_i < s_k$, whose source and target read the words $\altw{s_i}{s_k}m$
  and $\altw{s_k}{s_i}m$, $m = \cox(i,k)$, since $\varepsilon\cdot\altw stj$ is the unique element
  of $P(\{x_{j+1}\})$.
\end{proof}

\section{Application to precubical sets}\label{sec:cubes}

For a shape $a$ on $N$ events let $B_a := [N-1] \setminus \cut(a)$, the set of positions of its block
boundaries.

\begin{proposition}\label{prop:cubes}
  Sending $a$ to $B_a$, and a map of wedges $f$ to $f_*$, is an isomorphism from the full
  subcategory of $\Ch(Z)$ on the shapes with $N$ events onto $\EN$.
\end{proposition}
\begin{proof}
  On objects this is \texttt{lem:cuts}. On morphisms, $f \mapsto f_*$ is a functor because $\Ev$ is,
  and it is injective by \texttt{prop:coords}(i). For $f : \bigvee a \to \bigvee b$,
  \texttt{lem:young} gives $\cut(a) \subseteq \cut(b)$, that is $B_a \supseteq B_b$, and $\blk_b
  \circ f_* = \blk_b$, which is (x) for $B_b$; and \texttt{prop:coords}(ii) is (y) for $B_a$.
  TODO: surjectivity, that is, for $\phi \in \EN(B_a,B_b)$, a map of wedges with $f_* = \phi$ (Lean:
  \texttt{exists\_crossPerm}).
\end{proof}

Refinements preserve the number of events, so $\Ch(K)$, $\Br(K)$ and $\Art(K)$ split as disjoint
unions over $N$, and we work at one $N$ at a time. For a bipointed precubical set $K$, let $N(K)$ be
the presheaf $a \mapsto \BPSet(\bigvee a, K)$ on $\Ch(Z)$, read on $\EN$ through \Cref{prop:cubes}.
By the definition of $\Ch(K)$, $\Ch(K) = \int N(K)$; the order-preserving refinements are the
morphisms over $1$, so $\Br(K) = \Br(N(K))$; and $\Art(K) = \Art(N(K))$, with $\src$ and $\tgt$ given
by \Cref{def:art}. Theorem~1.1 of \texttt{mainv2} follows: (i) is \Cref{thm:A} at $P = N(K)$; (ii) is
its naturality clause, $N$ being a functor; (iii) is \Cref{cor:artin}, since $N(Z) = 1$.

\printbibliography

\end{document}
